\documentclass[11pt]{article}
\usepackage[margin=1in]{geometry}
\usepackage{booktabs}
\usepackage{longtable}
\usepackage{hyperref}
\usepackage{xcolor}
\usepackage{titlesec}
\usepackage{enumitem}
\hypersetup{colorlinks=true, urlcolor=blue, linkcolor=black}

\title{Replication Report: Zheng et al.\ (2017)\\
\large ``Complete genome sequencing and genomic characterization of two \textit{Escherichia coli} strains co-producing MCR-1 and NDM-1 from bloodstream infection''}
\author{Ollie (OpenClaw AI) --- BV-BRC Replication Project}
\date{2026-07-02}

\begin{document}
\maketitle

\begin{center}
\textbf{Paper:} Zheng B et al., \textit{Scientific Reports} 7:17885 (2017).\\
\textbf{DOI:} 10.1038/s41598-017-18273-2 \quad \textbf{PMID:} 29263349 \quad \textbf{PMCID:} PMC5738369\\
\textbf{Set:} BVBRC-58 \quad \textbf{Wave:} night push 2026-07-01/02\\
\textbf{Verdict:} \textbf{\textcolor{orange}{PARTIAL REPLICATION (strong)}}
\end{center}

\section{Paper summary}

The authors previously reported two \textit{E.\ coli} bloodstream isolates --- \textbf{EC1002} (ST405) and \textbf{EC2474} (ST131) --- that co-harbour the plasmid-borne colistin-resistance gene \textbf{mcr-1} and the carbapenemase \textbf{bla$_{\text{NDM-1}}$}. They complete-sequenced both isolates (PacBio RS II + Illumina HiSeq), producing two closed chromosomes plus 4 (EC1002) and 3 (EC2474) circular plasmids. Central findings:

\begin{itemize}[nosep]
  \item The two resistance genes reside on \textbf{different plasmids} in each strain (mcr-1 and bla$_{\text{NDM-1}}$ are never co-located).
  \item Diverse plasmid replicon backbones (mcr-1: IncI2 vs IncHI2; bla$_{\text{NDM-1}}$: IncA/C2 vs IncF).
  \item Two distinct mcr-1 mobilization contexts (\textit{nikA-nikB-mcr-1-hp} vs ISApl1-\textit{mcr-1}) and two distinct bla$_{\text{NDM-1}}$ contexts.
\end{itemize}

\section{Claims tested}

\begin{longtable}{clll}
\toprule
\# & Claim & Type & Tested? \\
\midrule
C1 & GenBank CP021202--CP021210 (2 chromosomes + 7 plasmids) & Data availability & \checkmark \\
C2 & Per-replicon genome statistics (sizes, GC) per Table 1 & Genome stats & \checkmark \\
C3 & EC1002 = ST405; EC2474 = ST131 & MLST & \checkmark \\
C4 & Per-plasmid acquired-resistance gene complement & Genomic (AMR) & \checkmark (partial) \\
C5 & Plasmid replicon types (IncI2/IncHI2/IncA/C2/IncF/\ldots) & Genomic (typing) & \checkmark \\
C6 & mcr-1 and bla$_{\text{NDM-1}}$ on \textbf{separate plasmids} & Central conclusion & \checkmark \\
\bottomrule
\end{longtable}

\section{Method}

All analyses were run on the \textbf{actual deposited GenBank sequences}, re-downloaded independently.

\begin{enumerate}[nosep]
  \item \textbf{Paper fetch:} Europe PMC full-text XML for PMC5738369 (\texttt{work/paper\_fulltext.xml}, 82~KB). Table~1 parsed.
  \item \textbf{Genome download:} NCBI \texttt{efetch.fcgi?db=nuccore\&id=<ACC>\&rettype=fasta} for all 9 accessions $\to$ \texttt{work/genomes/*.fasta}.
  \item \textbf{Genome statistics:} local venv, Biopython 1.87; \texttt{work/genome\_stats.py} computes length + GC\% per replicon $\to$ \texttt{evidence/evidence\_genome\_stats.json}.
  \item \textbf{MLST:} mlst 2.35.0, PubMLST \texttt{ecoli\_achtman\_4}: \\
    \texttt{mlst --scheme ecoli\_achtman\_4 strains/EC1002.fasta strains/EC2474.fasta} $\to$ \texttt{evidence/mlst\_results.tsv}.
  \item \textbf{Acquired resistance:} AMRFinderPlus 3.12.8 (DB 2024-07-22.1): \\
    \texttt{amrfinder -n strains/<S>.fasta --organism Escherichia --plus -d <db>} $\to$ \texttt{evidence/EC*\_amr.tsv}.
  \item \textbf{Plasmid replicon typing:} PlasmidFinder DB (\texttt{enterobacteriales.fsa}, 159 references) $\to$ \texttt{makeblastdb} $\to$ \texttt{blastn -perc\_identity 95}, coverage $\geq$60\% $\to$ \texttt{evidence/plasmidfinder\_results.tsv}.
  \item \textbf{Scoring:} LLM judge (Argo \texttt{gpt-5.2}, free), per-claim verdicts + coverage/agreement $\to$ \texttt{evidence/llm\_judge\_gpt52.md}.
\end{enumerate}

\section{Results vs paper}

\subsection{C1/C2 --- Accessions and genome statistics}

\begin{longtable}{lllrrrrr}
\toprule
Accession & Replicon & Paper bp & Obs bp & $\Delta$bp & Paper GC & Obs GC \\
\midrule
CP021202 & EC1002 chr        & 5{,}177{,}501 & 5{,}177{,}498 & $-3$ & 50.1 & 50.61 \\
CP021203 & pEC1002-1 IncFII  &   183{,}509  &   183{,}508  & $-1$ & 50.0 & 49.96 \\
CP021205 & pEC1002-MCR IncI2 &    63{,}392  &    63{,}392  & $ 0$ & 43.0 & 43.01 \\
CP021206 & pEC1002-NDM IncA/C2 & 111{,}688  &   111{,}688  & $ 0$ & 52.3 & 52.33 \\
CP021204 & pEC1002-4 IncFIB  &    92{,}439  &    92{,}438  & $-1$ & 50.0 & 50.30 \\
CP021207 & EC2474 chr        & 5{,}013{,}813 & 5{,}013{,}813 & $ 0$ & 50.6 & 50.61 \\
CP021209 & pEC2474-MCR IncHI2 &  223{,}982  &   223{,}982  & $ 0$ & 45.8 & 45.83 \\
CP021210 & pEC2474-NDM IncFII &   75{,}553  &    75{,}553  & $ 0$ & 50.8 & 50.78 \\
CP021208 & pEC2474-3 IncI1   &    86{,}717  &    86{,}725  & $+8$ & 49.5 & 49.55 \\
\bottomrule
\end{longtable}

All 9 replicons present; lengths match to 0--8~bp, GC to $\leq$0.5\%. \textbf{C1, C2 reproduced.}

\subsection{C3 --- MLST}

\begin{longtable}{lllll}
\toprule
Strain & Paper & Replication & Allelic profile \\
\midrule
EC1002 & ST405 & \textbf{ST405} \checkmark & adk35 fumC37 gyrB29 icd25 mdh4 purA5 recA73 \\
EC2474 & ST131 & \textbf{ST131} \checkmark & adk53 fumC40 gyrB47 icd13 mdh36 purA28 recA29 \\
\bottomrule
\end{longtable}

\textbf{C3 reproduced (exact).}

\subsection{C4 --- Per-plasmid acquired-resistance genes}

\textbf{EC1002}
\begin{longtable}{lllc}
\toprule
Replicon & Paper (ResFinder 2.1) & AMRFinderPlus 2024 & Assessment \\
\midrule
chr CP021202 & blaCTX-M-15, oqxB, tetB & blaCTX-M-15, tet(B), blaEC, QRDR muts & \checkmark core \\
pEC1002-1 CP021203 & blaCTX-M-15, sul, mph, aac(3)-Ib, erm, aadA4, dfrA, arr & blaCTX-M-15, aac(3)-IIe, aadA5, dfrA17, erm(B), mph(A) & \checkmark mostly \\
pEC1002-4 CP021204 & blaTEM & blaTEM-1 & \checkmark \\
\textbf{pEC1002-MCR CP021205} & \textbf{mcr-1 (only)} & \textbf{mcr-1.1 (only)} & \checkmark exact \\
\textbf{pEC1002-NDM CP021206} & blaNDM-1, blaCTX-M-14, blaTEM, sul1, mph, aac(6')-Ib, rmtC, arr & blaNDM-1, blaCTX-M-14, blaTEM-1, sul1, mph(A), aac(6')-Ib3, rmtC, ble, qacE$\Delta$1, aac(3)-IId & \checkmark strong \\
\bottomrule
\end{longtable}

\textbf{EC2474}
\begin{longtable}{lllc}
\toprule
Replicon & Paper & AMRFinderPlus 2024 & Assessment \\
\midrule
chr CP021207 & blaCTX-M-55 & blaCTX-M-55, blaEC & \checkmark \\
pEC2474-3 CP021208 & blaCTX-M-55 & blaCTX-M-55 & \checkmark \\
\textbf{pEC2474-MCR CP021209} & mcr-1, blaCTX-M-14, floR, aph4, sul2, aac(3)-IVa, fosA14 & mcr-1.1, blaCTX-M-14, floR, aph(4)-Ia, sul2, aac(3)-IVa, fosA3, ter* & \checkmark strong \\
\textbf{pEC2474-NDM CP021210} & blaNDM-1, aph & blaNDM-1, aph(3')-VI, ble & \checkmark \\
\bottomrule
\end{longtable}

Core AMR content reproduced across all plasmids; discrepancies are database-version artifacts. \textbf{C4 partially reproduced.}

\subsection{C5 --- Plasmid replicon typing}

\begin{longtable}{lllc}
\toprule
Plasmid & Paper & Replication & Match \\
\midrule
pEC1002-1 (CP021203) & IncFII & IncFII (100\%) & \checkmark \\
pEC1002-4 (CP021204) & IncFIB & IncFIB (98.8\%) + IncFIA & \checkmark \\
pEC1002-MCR (CP021205) & IncI2 & IncI2 (100\%) & \checkmark \\
pEC1002-NDM (CP021206) & IncA/C2 & IncC (100\%) & \checkmark (renamed) \\
pEC2474-3 (CP021208) & IncI1 & IncI1-I(Alpha) (100\%) & \checkmark \\
pEC2474-MCR (CP021209) & IncHI2 & IncHI2 (100\%) & \checkmark \\
pEC2474-NDM (CP021210) & IncF & IncFII (100\%) & \checkmark (refined) \\
\bottomrule
\end{longtable}

All 7 replicon types match. \textbf{C5 reproduced.}

\subsection{C6 --- mcr-1 and bla$_{\text{NDM-1}}$ on separate plasmids}

\begin{itemize}[nosep]
  \item EC1002: mcr-1.1 on \textbf{CP021205} (pEC1002-MCR, IncI2); bla$_{\text{NDM-1}}$ on \textbf{CP021206} (pEC1002-NDM, IncA/C2). Distinct plasmids. \checkmark
  \item EC2474: mcr-1.1 on \textbf{CP021209} (pEC2474-MCR, IncHI2); bla$_{\text{NDM-1}}$ on \textbf{CP021210} (pEC2474-NDM, IncF/FII). Distinct plasmids. \checkmark
  \item AMRFinderPlus additionally places bla$_{\text{NDM-1}}$ co-located with \textbf{rmtC + ble} on CP021206, consistent with the paper's \textit{rmtC-ISKpn14-bla}$_{\text{NDM-1}}$\textit{-bleMBL} context.
\end{itemize}

\textbf{C6 reproduced.} The paper's headline claim holds independently.

\section{LLM-judge scoring (Argo gpt-5.2, free)}

\begin{itemize}[nosep]
  \item \textbf{Coverage:} 5/6 = \textbf{83.3\%} of testable claims reproduced (C4 partial).
  \item \textbf{Agreement:} \textbf{$\sim$85--90\%}; agreement drops only on the gene-by-gene AMR inventory.
  \item \textbf{Canonical verdict:} \textbf{PARTIAL.}
\end{itemize}

\section{Genuine Critique}

\subsection{What genuinely replicated}

The \emph{sequence-substrate} half of this paper replicates cleanly and independently. Every claim that can be tested by re-downloading the 9 GenBank accessions and re-running standard genomic tools reproduces:

\begin{enumerate}[nosep]
  \item All 9 accessions (CP021202--CP021210) are live and retrievable via free NCBI efetch (C1).
  \item Per-replicon lengths match to \textbf{0--8~bp} and GC to \textbf{$\leq$0.5\%} (C2). The one apparent 0.5\% GC gap on the chromosome reflects paper-side rounding vs strict per-sequence computation, not a discrepancy.
  \item MLST is \textbf{exact}: EC1002 = ST405, EC2474 = ST131, allelic profiles reproduced digit-for-digit (C3).
  \item Plasmid replicon typing is \textbf{7/7} once one accepts standard renaming (IncA/C2 $\to$ IncC) and refinement (IncF $\to$ IncFII) (C5).
  \item The headline biological conclusion --- \textbf{mcr-1 and bla$_{\text{NDM-1}}$ sit on separate plasmids in both strains} --- reproduces exactly on the deposited sequences (C6).
\end{enumerate}

\subsection{What is only partial (and why)}

The per-plasmid \emph{acquired-resistance gene complement} (C4) reproduces at the \textbf{core-gene level} but drifts at the allele/nomenclature level. This is a database-version artifact, not a scientific disagreement:

\begin{itemize}[nosep]
  \item The paper used \textbf{ResFinder 2.1 (2017)}; this replication uses \textbf{AMRFinderPlus 3.12.8 with the 2024-07-22.1 database}. Seven years of curation separate them.
  \item AMRFinderPlus does not natively call \texttt{oqxB} or \texttt{arr} in this pipeline scope; those absences are calling-scope differences, not sequence-content differences.
  \item Newer databases annotate loci that were unnamed or unlisted in 2017 (e.g., \texttt{ble}, \texttt{qacE$\Delta$1}, tellurium-resistance clusters).
  \item Allele refinement (\texttt{mcr-1} $\to$ \texttt{mcr-1.1}, \texttt{fosA14} $\to$ \texttt{fosA3}, \texttt{aac(6')-Ib} $\to$ \texttt{aac(6')-Ib3}) is expected as reference variants have expanded.
\end{itemize}

None of the C4 differences challenge the paper's claims; they reflect the field's normal database drift.

\subsection{What was NOT replicated (limitations)}

\begin{itemize}[nosep]
  \item \textbf{Raw-read re-assembly:} PacBio RS II + Illumina HiSeq reads were not re-processed. Replication used the authors' \emph{closed} deposited sequences, which is appropriate for verifying genome-characterization claims but does not independently validate the assembly itself.
  \item \textbf{Comparative plasmid figures:} BRIG / Easyfig alignments (paper Figs.\ 2--3) were not rebuilt; genetic context of bla$_{\text{NDM-1}}$ was confirmed only at the co-localization level (bla$_{\text{NDM-1}}$ + \texttt{rmtC} + \texttt{ble} on one contig, consistent with paper's ISKpn14-based context but not aligned base-by-base).
  \item \textbf{Mobilization-element structure} (\texttt{nikA-nikB-mcr-1-hp} vs ISApl1-\texttt{mcr-1}) was not directly verified via IS/transposase annotation.
  \item \textbf{Conjugation/mobility assays} (in vitro transfer of mcr-1 or bla$_{\text{NDM-1}}$ plasmids) are wet-lab and out of scope.
\end{itemize}

\subsection{Verdict rationale}

The paper's \emph{central genomic conclusion} --- co-carriage of mcr-1 and bla$_{\text{NDM-1}}$ on separate, diverse plasmid backbones in bloodstream \textit{E.\ coli} --- is independently supported by re-analysis of the actual GenBank record. Genome stats, MLST, and replicon typing are essentially exact. Only per-gene AMR nomenclature drifts, for well-understood reasons. The overall replication is therefore \textbf{PARTIAL (strong)}: not FULL (because C4 has partial agreement and mobilization/context figures were not rebuilt), but far above STRUCTURAL --- the paper's biology holds up.

\section{Verdict}

\textbf{PARTIAL (strong).} Coverage 5/6 = 83.3\%; agreement $\sim$85--90\%.

\end{document}
